\chapter{Ownership}\label{ch-ownership}

A collector has two costs that firmware notices. It pauses the
program when it runs, and memory comes back at no particular moment. A
buffer filled a thousand times a second inside a control loop wants
neither. So a value in Zeno can be marked \verb|owned|. An owned value
has exactly one owner, and is freed the moment its owner's block ends,
at a line you can point to.

Shared values stay the default. Nothing changes for a program that
does not say \verb|owned|.

\section{Three words}

\begin{zeno}
owned class Packet {
    id: int
    data: owned list[int]
}

function make(id: int) -> Packet {
    let data: owned list[int] = [id, id * 2, id * 3]
    return Packet(id, data)       # `data` moves into the packet
}

function size(packet: Packet) -> int {      # looks
    return len(packet.data)
}

function send(packet: take Packet) -> int {  # takes
    return packet.id
}

let packet = make(7)
print(size(packet), packet.data)
print(send(packet))               # `packet` moves into `send`
\end{zeno}

\begin{itemize}
\item \verb|owned class|: every object of the class has exactly one
  owner. It is decided once, where the class is declared.
\item \verb|owned list[T]|: a list with exactly one owner.
\item \verb|take| on a parameter: the function takes the value over. A
  parameter without it only \emph{looks}: the function may read the
  value and change what is inside it, and may not keep it.
\end{itemize}

There are no lifetimes to write, and no \verb|&|.

\section{The rules}

\paragraph{1. One owner.} A variable, a parameter, or a field of an
owned object holds the value. Nothing else does.

\paragraph{2. Storing it gives it away.} Assigning it to another
variable, passing it to a \verb|take| parameter, putting it in a field,
or returning it moves it. The old name holds nothing afterwards, and
using it is an error that says where it went:

\begin{zeno}[fails=yes]
owned class Packet {
    id: int
}
function send(packet: take Packet) -> int {
    return packet.id
}
let packet = Packet(1)
print(send(packet))
print(packet.id)
\end{zeno}

\paragraph{3. Looking is free.} Reading a field, calling a method,
printing, indexing, comparing and looping over it move nothing.

\paragraph{4. A parameter looks, unless it says take.} What is only
looked at cannot be kept, returned or passed on to something that keeps
it, so it can never outlive the call:

\begin{zeno}[fails=yes]
owned class Packet {
    id: int
}
function keep(packet: Packet) -> Packet {
    return packet
}
\end{zeno}

\paragraph{5. It is freed where its owner's block ends:} at the
closing \verb|}|, or at a \verb|return|, \verb|break| or
\verb|continue| that leaves the block early. A \verb|var| given a new
value frees the old one at that line.

\paragraph{6. Owned may hold shared; shared may not hold owned.} A
shared list can be reached from many places, and if it held an owned
value, each of them would think the value was still there after its
owner had freed it. So a list of owned objects must itself be owned:

\begin{zeno}[fails=yes]
owned class Packet {
    id: int
}
let queue: list[Packet] = []
\end{zeno}

\begin{zeno}
owned class Packet {
    id: int
}
var queue: owned list[Packet] = []
queue.append(Packet(1))
queue.append(Packet(2))
let first = queue.pop(0)          # taken out: `first` owns it now
print(first, queue)
\end{zeno}

An owned list owns its items, and frees them with itself. A
\verb|Packet or none| can own a packet or nothing.

\section{Branches and loops}

The checker follows each path. After an \verb|if|, a value is gone if
any branch gave it away, unless that branch left with \verb|return|:

\begin{zenofrag}
if not ok {
    discard(frame)
    return
}
send(frame)          # fine: only reached with the frame still here
\end{zenofrag}

A value given away on one trip round a loop would not be there on the
next, so an owner from outside a loop cannot be given away inside it,
except just before a \verb|break|:

\begin{zeno}[fails=yes]
owned class Packet {
    id: int
}
function send(packet: take Packet) {
    print("sent", packet.id)
}
let packet = Packet(1)
for turn in range(3) {
    send(packet)
}
\end{zeno}

In a \verb|try|, an error may leave the body after any line, so the
\verb|except| starts with everything the body gives away already gone.

\section{Seeing the frees}

This program from the \verb|examples| folder makes 120 frames of 600
samples each, sends some and skips others. Every list and object in it
is owned:

\begin{zeno}[src=examples/buffer.zeno, cmd=zeno build --memory buffer.zeno > /dev/null ;; ./buffer, file=buffer.zeno]
\end{zeno}

It made megabytes of buffers, never held more than a few kilobytes,
and the collector never ran: every block was freed at a line, the
moment its owner was done.

\begin{why}
Rust showed that one owner, checked by the compiler, gives memory back
at exact moments with no collector and no mistakes. It also showed that
borrowing everywhere, with lifetimes written in the program, is the
hardest thing in Rust to learn. Zeno takes the first half. A parameter
looks, which needs nothing written, and only keeping needs a word. No
lifetimes are needed because a look can never outlive the call.
Owned values are a choice made where it pays, and shared values stay
the easy default. (Decisions 5.8 and 5.15, and the audit, which made
looking the default.)
\end{why}

\section{Taking a value out of a field}

Reading a field only looks, so an object could not hand on what one of
its fields holds: giving it away would leave nothing in the field.
\verb|take| in front of a field or an item gives what is there and
leaves \verb|none| in its place, so the field's type must say it may
be \verb|none|:

\begin{zeno}
owned class Buffer {
    bytes: list[int]
}

owned class Uart {
    filling: Buffer or none

    function hand_on(self) -> Buffer or none {
        return take self.filling
    }
}

function send(buffer: take Buffer) -> int {
    return len(buffer.bytes)
}

let uart = Uart(Buffer([1, 2, 3]))
let full = uart.hand_on()
print(uart.filling is none)
if full is not none {
    print(send(full))
}
\end{zeno}

\begin{why}
A list could always give an item away with \verb|pop()|, and a field
had nothing like it. \verb|take| is the word a parameter already uses
for the same act, and writing it at the place makes the hole visible:
the type says \verb|or none|, and the next reader of the field must
test it. (Decision 56.10; new in 0.8.)
\end{why}

\section{What owned values cannot do}

An owned value cannot be replaced through a look, or taken out of a
field whose type does not say \verb|or none| (that would leave a hole
the type does not admit to). Each is an error that says what to do
instead.

\begin{summary}
\verb|owned class| and \verb|owned list[T]| have one owner, and are
freed where the owner's block ends. Storing gives a value away; looking
is free. A parameter looks unless it says \verb|take|. Owned may hold
shared, never the other way. The checker follows branches, loops and
\verb|try|, and says where a value went when it is used after.
\end{summary}

\begin{tryit}
\begin{enumerate}
\item Change \verb|size| above to return the packet's list. Read the
  message, and say in your own words why it is right.
\item Write a queue of owned packets with \verb|put| and \verb|get|
  methods, where \verb|get| gives \verb|Packet or none|.
\item In \verb|buffer.zeno|, remove \verb|take| from \verb|send|. What
  changes in the report?
\item Give an owned value away inside a loop and then \verb|break| on
  the next line. Why does the checker accept that?
\end{enumerate}
\end{tryit}
